%-------------------------------------------------------------------------------
%	ABSTRACT PAGE
%-------------------------------------------------------------------------------

\addchaptertocentry{Хураангуй} % Хураангуйг гарчигт нэмэх

\begin{center}
{\scshape\Large \univname\par} % Их сургуулийн нэр
{\scshape\large \facname\par}\vspace{0.5cm} % Их сургуулийн нэр
{\huge\textbf{{Хураангуй}} \par}
\bigskip
{\Large{\ttitle} \par} % Тезисийн нэр
\bigskip

{\normalsize \shortname \par}\addressname

% Зохиогчийн нэр

\end{center}

\textit{\textbf{Түлхүүр үгс: \keywordnames}}
\paragraph{}Монголын цахим орчин дахь сэтгэгдэл, санал бодол нь олон нийтийн хандлага, үйлчилгээний чанар болон байгууллагын хариуцлагыг үнэлэх чухал эх сурвалж болж байна. Эдгээр өгөгдөлд хувь хүн, бүлэг рүү чиглэсэн доромжлол, дарамт, заналхийлэл бүхий хортой сөрөг агуулга болон үйлчилгээний чанар, байгууллагын үйл ажиллагаа эсвэл системийн доголдолд чиглэж, сайжруулах санал агуулсан бүтээлч шүүмжлэл зэрэгцэн оршдог тул бүх сөрөг текстийг нэг ангилалд нэгтгэн авч үзэх нь хангалтгүй юм.

\paragraph{}Энэхүү асуудлыг шийдвэрлэхийн тулд Монголын цахим орчны кирилл текст дэх сэтгэгдлийг дөрвөн ангилалд (эерэг, саармаг, хортой сөрөг, бүтээлч шүүмжлэл) автоматаар хуваарилах эх хэлний боловсруулалтын (NLP) системийг боловсрууллаа. news.mn, gogo.mn болон Facebook зэрэг нийтэд нээлттэй эх сурвалжуудаас 10,000 сэтгэгдлийг цуглуулж, гурван бие даасан шошгологчийн тусламжтайгаар үүсгэсэн өгөгдлийн багцын нийцлийг Флейссийн каппа ($\kappa{=}0.72$) үзүүлэлтээр баталгаажуулсан. Кирилл тэмдэгтийн шалгуур, HTML тэмдэглэгээ болон URL холбоос хасалт, эможи орлуулалт зэрэг урьдчилсан боловсруулалтын дамжлагыг залгамал морфологийн онцлогт тохирсон SentencePiece токенчлолтой нэгтгэж, MN-BERT загварыг нарийвчлан тохируулан (fine-tuning) үгийн давтамж-баримтын ялгарлын жинд (TF-IDF) суурилсан логистик регресс, Наив Байес, тулгуур векторын машин (SVM) загваруудтай харьцуулан үнэлэв.

\paragraph{}Туршилтаар хоёр шатлалт MN-BERT нэгдсэн дамжлага эцсийн дөрвөн ангиллын түвшинд $78.94\%$ нарийвчлал, $0.7628$ макро-F1 үзүүлсэн бөгөөд шат дамжсан алдаа нь энэ зохион байгуулалтын гол хязгаарлалт болох нь ажиглагдав. Харин фокус алдагдлын функц (Focal Loss), ангиллын жин, тэнцвэржүүлсэн түүвэрлэлт болон косинус хэлбэрийн сургалтын халаалтыг (Cosine Warmup) хослуулсан нэг шатлалт MN-BERT загвар $\mathbf{83.26\%}$ нарийвчлал, $\mathbf{0.8062}$ макро-F1 оноо гаргав. Энэхүү нэг шатлалт бүтэц нь уламжлалт машин сургалтын суурь загварын хамгийн өндөр нарийвчлал ($66.91\%$) болон өмнөх MN-BERT ажлын нарийвчлалын үзүүлэлтээс ($82.70\%$)~\cite{jargalmaa2025} бага хэмжээгээр өндөр гарч, хортой агуулга болон бүтээлч шүүмжлэлийг Монгол кирилл текстэд автоматаар ялган таних боломжтойг харууллаа.
